%! AP Statistics — Unit 4, Lesson 4.10 — Warm-Up KEY
\documentclass[10pt]{article}
\usepackage{apstats-article}
\usepackage{apstats-key}

\begin{document}

\pageheader{Unit 4, Lesson 4.10}{Warm-Up --- Answer Key}
\namedateperiod

\vspace{0.1in}

% ── SECTION 1: Spiral Review ──────────────────────────────────────────────────
{\large\bfseries\color{navy} Part 1 \quad Spiral Review}

\vspace{0.08in}

\begin{enumerate}

    \item A discrete random variable $X$ has the following partial probability
    distribution. Find the missing probability and the mean $\mu_X$.

    \vspace{8pt}
    \renewcommand{\arraystretch}{1.8}
    \begin{tabularx}{\linewidth}{@{} >{\bfseries}p{2cm} X X X X @{}}
    $x$ & 0 & 1 & 2 & 3 \\
    \hline
    $P(X=x)$ & 0.10 & 0.35 & \ans{0.30} & 0.25 \\
    \end{tabularx}

    \vspace{8pt}
    Missing probability $= $ \ans{0.30} \hspace{1cm} $\mu_X = $ \ans{0(0.10)+1(0.35)+2(0.30)+3(0.25) = 1.70}

    \vspace{0.1in}

    \item Evaluate each combination by hand. Show your work using $\binom{n}{k} = \dfrac{n!}{k!(n-k)!}$.

    \vspace{8pt}
    \begin{tabularx}{\linewidth}{@{} X X X @{}}
    (a) $\displaystyle\binom{5}{2} = $ \ans{10} &
    (b) $\displaystyle\binom{6}{0} = $ \ans{1} &
    (c) $\displaystyle\binom{4}{4} = $ \ans{1} \\[10pt]
    \end{tabularx}

    \vspace{0.1in}

    \item Two events $A$ and $B$ are independent with $P(A) = 0.6$ and $P(B) = 0.4$.
    \begin{enumerate}[label=(\alph*), itemsep=3pt]
        \item Find $P(A \cap B)$.\\
        \ansline{$P(A)\cdot P(B) = 0.6 \times 0.4 = 0.24$}
        \item Find $P(\text{neither } A \text{ nor } B)$.\\
        \ansline{$(1-0.6)(1-0.4) = 0.4 \times 0.6 = 0.24$}
    \end{enumerate}

    \vspace{0.1in}

    \item A bag contains 4 red and 6 blue marbles. You draw one marble, note its color,
    and replace it before drawing again. You do this 3 times.
    \begin{enumerate}[label=(\alph*), itemsep=3pt]
        \item Are the three draws independent? Why?\\
        \ansline{Yes --- with replacement the bag is the same for every draw.}
        \item What is the probability that all three draws are red?\\
        \ansline{$(0.4)^3 = 0.064$}
    \end{enumerate}

\end{enumerate}

\vspace{0.08in}

% ── SECTION 2: Looking Ahead ──────────────────────────────────────────────────
{\large\bfseries\color{navy} Part 2 \quad Looking Ahead}

\vspace{0.08in}

\noindent Today we use the structure from Problem 4 --- repeated independent trials
with the same probability --- to build the \textbf{binomial distribution}.

\vspace{0.1in}

\begin{tcolorbox}[colback=greenbg, colframe=greenacc, boxrule=0.8pt, arc=2mm,
    left=4mm, right=4mm, top=3mm, bottom=3mm]
    \small
    \begin{itemize}[leftmargin=1.3em, itemsep=4pt]
  \item A binomial setting requires four conditions (BINS) to be verified from context.
  \item The binomial formula counts both the arrangements and the probabilities.
  \item Always check BINS before using any binomial formula or calculator command.
    \end{itemize}
\end{tcolorbox}

\vspace{0.08in}

\noindent\textcolor{slate}{\small Write one question you have about today's lesson:}\\[8pt]
\ansline{Student responses will vary.}\\[2pt]\ansline{}

\end{document}
